\begin{lemma}
Let $H$ be a finite indexed $k$-uniform hypergraph of maximum degree at
most $D$, and suppose that an edge index $e$ has exactly $k(D-1)$ distinct
neighbors. For $v\in e$ put $C_v=\{g\ne e:v\in H(g)\}$.
Then $|C_v|=D-1$ for every $v\in e$, these classes are pairwise disjoint,
and their union is precisely the neighbor set of $e$. Moreover,
$|H(g)\cap H(e)|\leq1$ for every $g\ne e$.
\end{lemma}
\begin{proof}
By \noderef{MaxDegreeAtMost}, removing the incident index $e$ leaves
$|C_v|\leq D-1$ at each $v\in e$. By \noderef{UniformHypergraph},
$|H(e)|=k$. Let $N$ be the neighbor set. Counting pairs $(v,g)$ with
$v\in H(e)\cap H(g)$ and $g\ne e$ gives
\[
k(D-1)=|N|\leq\sum_{g\in N}|H(g)\cap H(e)|
=\sum_{v\in H(e)}|C_v|\leq k(D-1).
\]
Each term in the first sum is at least one, so equality forces each to
equal one. Similarly, each term in the last sum is at most $D-1$, so
equality forces each to equal $D-1$. An index outside $N\cup\{e\}$
has empty intersection with $H(e)$, proving the asserted bound for all
$g\ne e$. If an index lay in two classes with distinct labels, its
intersection with $H(e)$ would have at least two vertices, a contradiction.
Finally, membership in a class is exactly the existence of a vertex in
$H(g)\cap H(e)$ for an index $g\ne e$, proving the union identity.
All counts are of indices, so this argument allows parallel hyperedges.
\end{proof}
